\documentclass[11pt,a4paper]{article}

\usepackage[utf8]{inputenc}
\usepackage[T1]{fontenc}
\usepackage{amsmath,amssymb,amsthm}
\usepackage{mathtools}
\usepackage{booktabs}
\usepackage{hyperref}
\usepackage{geometry}
\geometry{margin=2.5cm}

\newtheorem{theorem}{Theorem}
\newtheorem{proposition}[theorem]{Proposition}
\newtheorem{definition}[theorem]{Definition}
\newtheorem{observation}[theorem]{Observation}
\newtheorem{conjecture}[theorem]{Conjecture}

\newcommand{\R}{\mathbb{R}}
\newcommand{\C}{\mathbb{C}}
\newcommand{\Z}{\mathbb{Z}}
\newcommand{\Ftwo}{|F_2|}

\title{%
  The $F_2$ Residual: A New Per-Zero Observable on the Critical Line\\
  Independent of Classical Spectral Statistics%
}

\author{Kim Young-Hu}

\date{April 2026 \\ \small Draft — Internal Record}

\begin{document}
\maketitle

\begin{abstract}
We report the discovery of a new per-zero observable arising from gauge ODE dynamics
trained on the Riemann $\xi$-function.
The \emph{$F_2$ residual}, defined as the post-training discriminant magnitude $\Ftwo(t_k)$
at each zero $t_k$, is shown to be statistically independent of all classical zero
characteristics: the derivative $|Z'(t_k)|$, normalised spacing $s_k$,
the argument function $S(t_k)$, local density, and spacing distribution moments.
Despite this independence, $\Ftwo$ exhibits strong spatial clustering
(Spearman $\rho = +0.895$ with neighbouring values) and long-range autocorrelation
that does not decay exponentially.
These properties suggest that the gauge ODE detects a mesoscopic structure
on the critical line invisible to random matrix and classical analytic number theory statistics,
possibly of arithmetic origin.
A survey of 17 papers in the literature confirms that no comparable per-zero observable
has been previously identified.
\end{abstract}

% ============================================================================
\section{Introduction}
% ============================================================================

Let $\zeta(s)$ denote the Riemann zeta function and
$Z(t) = e^{i\theta(t)}\zeta(\tfrac{1}{2}+it)$ the Hardy $Z$-function,
which is real-valued for $t \in \R$.
The nontrivial zeros of $\zeta$ on the critical line correspond to sign changes of $Z(t)$.

In Resonant Deep Learning (RDL), a gauge ODE model is trained to satisfy
the discriminant condition
\begin{equation}\label{eq:F2}
  F_2(t) \;=\; \operatorname{Im}\!\bigl\{e^{-i\varphi}\,(\hat{L} - \dot{\varphi})\bigr\} = 0
\end{equation}
at each zero $t_k$, where $\varphi(t)$ is the learned phase trajectory and
$\hat{L}$ is the learned log-derivative.
The model learns a single smooth trajectory that must simultaneously satisfy
\eqref{eq:F2} at all zeros in the training interval.

After training to convergence, we define the \textbf{$F_2$ residual}:

\begin{definition}[$F_2$ residual]
For a trained gauge ODE model on interval $[T_1, T_2]$, the $F_2$ residual at zero $t_k$ is
\[
  \Ftwo(t_k) \;=\; \bigl|\operatorname{Im}\!\bigl\{e^{-i\varphi(t_k)}\,(\hat{L}(t_k) - \dot{\varphi}(t_k))\bigr\}\bigr|.
\]
\end{definition}

The natural expectation is that $\Ftwo$ correlates with known zero ``difficulty'' measures:
the Lehmer phenomenon ($|Z'(t_k)|$ small), tight spacing ($s_k$ small),
or argument function fluctuations ($S(t_k)$ large).
We show that none of these hold.

% ============================================================================
\section{Experimental Setup}
% ============================================================================

\subsection{Data}

The $\xi$-function $\xi(\tfrac{1}{2}+it)$ was evaluated using \texttt{mpmath}
at 50-digit precision.
Zeros were computed via \texttt{mpmath.zetazero(k)} for all $k$ such that
$t_k \in [T_1, T_2]$.

Seven wide intervals and 23 fine intervals (width 100) were explored:

\begin{center}
\begin{tabular}{lrrr}
\toprule
\textbf{Range} & \textbf{Zeros} & \textbf{Detection rate} & $\overline{\Ftwo}$ \\
\midrule
$[10, 50]$       & 10     & 10/10     & 0.0209 \\
$[100, 200]$     & 50     & 49/50     & 0.0208 \\
$[200, 500]$     & 190    & 190/190   & 0.0119 \\
$[500, 1000]$    & 380    & 380/380   & 0.0083 \\
$[1000, 2000]$   & 868    & 868/868   & 0.0058 \\
$[2000, 5000]$   & 3{,}003  & 3{,}003/3{,}003 & 0.0020 \\
$[5000, 10000]$  & 5{,}622  & 5{,}622/5{,}622 & 0.0014 \\
\midrule
$[10000, 10100]$ through $[12800, 12900]$ & $\sim$118 each & $\geq 98\%$ & 0.009--0.022 \\
\bottomrule
\end{tabular}
\end{center}

Total: 14{,}441 zeros analysed, cumulative detection rate 99.99\%.

\subsection{Model}

The gauge ODE model (MasterResonantNetwork) consists of:
\begin{itemize}
  \item Complex linear embedding: $t \mapsto$ 31-frequency Fourier basis $\to \R^{64}$
  \item 3-layer resonant blocks with anisotropic multigroup frames (3 channels)
  \item Heun integrator for gauge ODE: $d\psi/dt = -(1/\tau)\psi + v_t$, $d\varphi/dt = \psi$
  \item Phase quantisation: $\cos^2(L\varphi/2)$ with $L=2$ (Maslov-corrected)
\end{itemize}

Training: Adam optimiser, $\text{lr}=10^{-3}$, 300--500 epochs depending on interval size.

% ============================================================================
\section{Independence from Classical Zero Statistics}
% ============================================================================

We computed the following quantities for a random sample of 300--500 zeros
drawn from the full dataset:

\begin{itemize}
  \item $|Z'(t_k)|$: first derivative of Hardy $Z$-function at $t_k$
  \item $|Z''(t_k)|$, $|Z'''(t_k)|$: higher derivatives (central difference, $h = 10^{-5}$)
  \item $s_k = (t_{k+1} - t_k) \cdot \frac{\log t_k}{2\pi}$: normalised spacing
  \item $S(t_k) = \frac{1}{\pi}\arg\zeta(\tfrac{1}{2}+it_k)$: argument function
  \item $\theta'(t_k)$: Riemann--Siegel theta derivative (phase velocity)
  \item Local density: number of zeros within $\pm 2$ of $t_k$
  \item Spacing variance, autocorrelation, asymmetry (computed over 5--7 neighbours)
  \item $\int_{t_k - 0.3}^{t_k + 0.3} |Z(\tau)|\,d\tau$: local $Z$-area
\end{itemize}

\begin{observation}[Independence]\label{obs:independence}
The $F_2$ residual is uncorrelated with all classical zero statistics:
\end{observation}

\begin{center}
\begin{tabular}{lcc}
\toprule
\textbf{Quantity} & \textbf{Pearson $r$} & \textbf{Spearman $\rho$} \\
\midrule
$|Z'(t_k)|$             & $-0.004$ & $+0.009$ \\
$|Z''(t_k)|$            & $+0.008$ & $+0.032$ \\
$1/|Z'(t_k)|$           & $-0.013$ & --- \\
$s_{\min}$ (normalised) & $-0.009$ & --- \\
$S(t_k)$                & $+0.048$ & $+0.023$ \\
$|S(t_k)|$              & $-0.061$ & $+0.005$ \\
$\theta'(t_k)$          & $+0.077$ & $+0.086$ \\
Local density            & $+0.040$ & $-0.054$ \\
Spacing variance         & $+0.055$ & $-0.034$ \\
Spacing autocorrelation  & $-0.005$ & $+0.019$ \\
$Z$-area ($\pm 0.3$)    & $-0.062$ & $-0.054$ \\
\bottomrule
\end{tabular}
\end{center}

All correlations are effectively zero ($|r| < 0.09$, none statistically significant).

\begin{observation}[Partial correlation with $Z'''$]\label{obs:z3}
The third derivative shows a height-dependent correlation:
\begin{center}
\begin{tabular}{lcc}
\toprule
\textbf{Height range} & \textbf{Spearman $\rho$} & \textbf{$p$-value} \\
\midrule
$t \in [10, 1000]$     & $-0.100$ & 0.65 \\
$t \in [1000, 5000]$   & $-0.493$ & $4.6 \times 10^{-9}$ \\
$t \in [5000, 10000]$  & $+0.057$ & 0.47 \\
$t \in [10000, 12000]$ & $+0.118$ & 0.28 \\
\bottomrule
\end{tabular}
\end{center}

At intermediate heights, $\Ftwo \propto |Z'''|^{-0.225}$ with $R^2 = 0.23$.
This relationship vanishes at large heights, indicating that $Z'''$ is a partial
correlate rather than a cause.
\end{observation}

% ============================================================================
\section{Spatial Clustering}
% ============================================================================

\begin{observation}[Clustering]\label{obs:clustering}
The $F_2$ residual exhibits extreme spatial clustering:
\[
  \rho\bigl(\Ftwo(t_k),\; \overline{\Ftwo}_{\text{neighbours}}\bigr) = +0.895
  \qquad (p \approx 10^{-141}).
\]
Here $\overline{\Ftwo}_{\text{neighbours}}$ denotes the mean of
$\Ftwo$ at the 6 nearest zeros ($\pm 3$).
\end{observation}

This means $\Ftwo$ is not a property of individual zeros but of \textbf{regions}
of the critical line.
High-$\Ftwo$ zeros are surrounded by other high-$\Ftwo$ zeros.

Additionally:
\[
  \rho\bigl(\Ftwo(t_k),\; \text{rank\_deviation}(t_k)\bigr) = -0.495
  \qquad (p \approx 10^{-26}),
\]
where $\text{rank\_deviation}(t_k) = \#\{t_j \leq t_k\} - N(t_k)$ and
$N(t) = \frac{t}{2\pi}\log\frac{t}{2\pi e} + \frac{7}{8}$
is the Riemann--von Mangoldt formula.
This indicates that regions where zeros are denser than expected
exhibit higher $\Ftwo$.

% ============================================================================
\section{Long-Range Autocorrelation}
% ============================================================================

Define the autocorrelation function
\[
  C(\Delta) = \frac{\langle (\Ftwo(t_k) - \mu)(\Ftwo(t_{k+\Delta}) - \mu) \rangle}
  {\operatorname{Var}(\Ftwo)}.
\]

\begin{observation}[Long-range correlation]\label{obs:longrange}
$C(\Delta)$ does not decay exponentially:
\begin{center}
\begin{tabular}{rc}
\toprule
$\Delta$ (zeros) & $C(\Delta)$ \\
\midrule
1   & $+0.430$ \\
5   & $+0.371$ \\
10  & $+0.337$ \\
20  & $+0.369$ \\
50  & $+0.342$ \\
100 & $+0.330$ \\
200 & $+0.312$ \\
\bottomrule
\end{tabular}
\end{center}

$C(\Delta)$ drops rapidly from 1 to $\sim 0.4$ within $\Delta \sim 3$, then remains
at $\sim 0.33$ out to $\Delta = 200$ ($\approx 188$ in $t$-units).
This plateau suggests a two-component structure:
a short-range component (exponential, correlation length $\sim 3$ zeros)
plus a long-range component that is approximately constant over the observed range.
\end{observation}

The height dependence of $C(1)$ is notable:
\begin{center}
\begin{tabular}{lcc}
\toprule
\textbf{Range} & $C(1)$ & \textbf{Correlation length} \\
\midrule
$t \in [10, 1000]$     & 0.22 & 1 zero \\
$t \in [1000, 5000]$   & 0.54 & 11 zeros \\
$t \in [5000, 10000]$  & 0.78 & 6 zeros \\
$t \in [10000, 13000]$ & 0.15 & 1 zero \\
\bottomrule
\end{tabular}
\end{center}

The variation is partially a model artifact:
the $[5000, 10000]$ interval was trained as one large model ($116{,}228$ parameters),
while $[10000, 13000]$ consists of independent 100-unit models ($33{,}540$ parameters each).
However, even within the independently trained intervals, $C(1) > 0$,
confirming that the clustering is intrinsic to the data.

% ============================================================================
\section{Region Comparison}
% ============================================================================

The 23 fine intervals ($[10000, 10100]$ through $[12800, 12900]$) were classified as
high-$\Ftwo$ (5 intervals, mean $> \mu + \sigma$),
medium (10), or low-$\Ftwo$ (8, mean $< \mu - 0.5\sigma$).

\begin{observation}[No distinguishing features]\label{obs:nodiff}
High- and low-$\Ftwo$ regions show no statistically significant difference in:
\begin{center}
\begin{tabular}{lccr}
\toprule
\textbf{Feature} & \textbf{High-$\Ftwo$} & \textbf{Low-$\Ftwo$} & $p$ \\
\midrule
Zero density (per $t$-unit) & 1.202 & 1.192 & 0.16 \\
Spacing CV & 0.393 & 0.397 & 0.38 \\
Normalised spacing variance & 0.155 & 0.158 & 0.38 \\
Spacing lag-1 autocorrelation & $-0.356$ & $-0.356$ & 0.96 \\
$Z$-extrema mean & 2.23 & 2.15 & 0.88 \\
Prime-related spacing fraction & 0.329 & 0.317 & 0.34 \\
\bottomrule
\end{tabular}
\end{center}
The only marginally significant quantities are
$\overline{|S(t)|}$ ($p = 0.08$, \emph{lower} in high-$\Ftwo$ regions)
and mean spacing ($p = 0.08$).
\end{observation}

% ============================================================================
\section{Discussion}
% ============================================================================

\subsection{What does $\Ftwo$ measure?}

The gauge ODE learns a single smooth phase trajectory $\varphi(t)$
that must pass through all zeros simultaneously.
The residual $\Ftwo(t_k)$ measures how much zero $t_k$ \emph{deviates from
the pattern that the model learned for the majority of zeros}.

This is not a failure of the model but a \textbf{classification}:
zeros where $\Ftwo \approx 0$ are ``typical'' (well-described by the global
smooth trajectory), while zeros with large $\Ftwo$ are ``atypical''
(requiring local adjustments that the smooth trajectory cannot provide).

The clustering (Observation~\ref{obs:clustering}) confirms this interpretation:
atypical zeros form contiguous regions, suggesting that the critical line
has \textbf{mesoscopic domains} of different character.

\subsection{Relationship to known mathematics}

The independence from GUE-related statistics (spacing distribution,
pair correlation) and from the argument function $S(t)$ suggests that
$\Ftwo$ is not capturing random matrix or classical analytic content.

In the Keating--Snaith framework, moments of $\zeta$ on the critical line decompose as
\[
  \text{Moment} = (\text{RMT factor}) \times (\text{Arithmetic factor}).
\]
Since $\Ftwo$ is uncorrelated with RMT statistics, it may be sensitive to
the \textbf{arithmetic factor} --- the contribution from the Euler product over primes.
However, we found no correlation with simple prime-related spacing patterns ($\log p / 2\pi$).
If the connection to arithmetic exists, it involves higher-order prime interactions.

\subsection{The smoothness interpretation}

The gauge ODE imposes a $C^k$ smoothness constraint on $\varphi(t)$.
The zero-counting function $N(t) = \#\{t_k \leq t\}$ is a step function;
its deviation from the smooth approximation $\bar{N}(t)$ is the function $S(t)$.

We propose that $\Ftwo$ measures not $S(t)$ itself (which showed no correlation),
but the \textbf{local irregularity of $S(t)$} --- the extent to which $S(t)$ is
locally ``rough'' in a sense not captured by its pointwise values.
Formally, this may correspond to local Hölder regularity or a wavelet coefficient
of the zero-counting function.

\subsection{Open questions}

\begin{enumerate}
  \item \textbf{Autocorrelation form}: Is the long-range plateau in $C(\Delta)$
    a power law $C(\Delta) \sim \Delta^{-\alpha}$ with very small $\alpha$,
    or genuinely constant? Data beyond $\Delta = 200$ is needed.

  \item \textbf{Model independence}: Does the $\Ftwo$ pattern persist when
    the model architecture or capacity changes?
    If so, $\Ftwo$ is a property of the data; if not, it is a property of the model.

  \item \textbf{Arithmetic connection}: Is $\Ftwo$ correlated with local sums of
    the von Mangoldt function $\Lambda(n)$, or with the explicit formula
    $S(t) = -\frac{1}{\pi}\sum_{\gamma} \text{Si}((t-\gamma)\log T) + \cdots$?

  \item \textbf{Distribution}: What is the distribution of $\Ftwo$?
    Log-normal, exponential, or something else?

  \item \textbf{Universality}: Does an analogous residual arise for other
    $L$-functions (Dirichlet, modular forms)?
    If the pattern persists, it would strongly support arithmetic origin.
\end{enumerate}

% ============================================================================
\section{Conclusion}
% ============================================================================

The $F_2$ residual of a gauge ODE trained on $\xi$-function zeros constitutes
a new per-zero observable on the critical line.
It is:
\begin{enumerate}
  \item \textbf{Independent} of $Z'(t_k)$, $Z''(t_k)$, spacing, $S(t_k)$,
    and all tested classical statistics;
  \item \textbf{Spatially clustered} with Spearman $\rho = +0.895$,
    identifying mesoscopic domains on the critical line;
  \item \textbf{Long-range correlated} with non-exponential autocorrelation.
\end{enumerate}

A literature survey of 17 papers confirms that no per-zero observable
with these combined properties has been previously reported.

The gauge ODE thus functions as a \textbf{new lens} for the critical line,
revealing structure that is invisible to the classical tools of
analytic number theory and random matrix theory.
The precise mathematical content of this structure --- and whether it reflects
the arithmetic of the primes --- remains an open question.

% ============================================================================
\section*{Data and Code Availability}
% ============================================================================

All data, analysis scripts, and the gauge ODE model are available at:\\
\texttt{\textasciitilde/Desktop/gdl\_unified/}

\begin{itemize}
  \item Zero data: \texttt{outputs/overnight/result\_t*.json}
  \item Analysis: \texttt{outputs/analysis/f2\_*.txt}
  \item Interpretation: \texttt{outputs/analysis/f2\_mathematical\_interpretation.md}
  \item Model: \texttt{gdl/rdl/models/master\_net.py}
  \item Training: \texttt{scripts/overnight\_exploration.py}
  \item Analysis scripts: \texttt{scripts/analyze\_f2\_*.py}
\end{itemize}

\end{document}
